% GENERATED FILE. Edit canonical lessons, then run npm run generate:books.
% canonical-source-hash: fnv1a64:a9fc153c408c09e3
% canonical-lessons: FR-C170-chateau, FR-C170-tour, FR-C170-monument, FR-C170-statue, FR-C170-fontaine

\chapter{Castle, Tower, Monument, Statue, Fountain}
\label{ch:fr-castle-tower-monument-statue-fountain}
\hlchaptermodality{170}

\begin{chapteropening}
\textbf{By the end of this chapter:} \emph{I can say a castle, a tower, a monument, a statue and a fountain.}

Complete the last lesson of chapter 170: I can say a castle, a tower, a monument, a statue and a fountain.
\end{chapteropening}

\section[le château]{le château — a castle}

\label{lesson:FR-C170-chateau}



\begin{quote}
Before the new one: say the French for a label, then the French for a guide, a guidebook.
\end{quote}

\subsection*{You'll want to know: le château}
\textbf{le château} — \textquotedblleft{}a castle\textquotedblright{}.

\begin{cousinweb}
\textbf{château} $\leftarrow$ Old French \emph{chastel} $\leftarrow$ Latin \emph{castellum}, \textquotedblleft{}a fort\textquotedblright{} (EtymDB 2.1) — inherited: \emph{ca-} became \emph{cha-}, and the circumflex marks the lost \emph{s}.
\end{cousinweb}

\begin{grammarlens}[title={What you've built}]
One more word for this chapter.
\end{grammarlens}

\subsection*{Your turn}
\begin{itemize}
\raggedright
  \item \emph{Say it:} \emph{le château}
  \item \emph{Say it:} \emph{le château}, once more
  \item \emph{Say it:} say \emph{le château}
  \item \emph{Recall:} say the French for a label, then the French for a guide, a guidebook, then say \emph{le château} again
\end{itemize}

\begin{tcolorbox}[breakable,colback=teal!4,colframe=teal!35!black,title={Before you move on}]
What does le château mean? (\textquotedblleft{}A castle\textquotedblright{}.) Say it once more.
\end{tcolorbox}
\section[la tour]{la tour — a tower}

\label{lesson:FR-C170-tour}



\begin{quote}
Before the new one: say the French for a guide, a guidebook, then the French for a castle.
\end{quote}

\subsection*{You'll want to know: la tour}
\textbf{la tour} — \textquotedblleft{}a tower\textquotedblright{}.

\begin{cousinweb}
\textbf{tour} $\leftarrow$ Old French \emph{tor} $\leftarrow$ Latin \emph{turris}, \textquotedblleft{}a tower\textquotedblright{}, from Greek — inherited. (EtymDB 2.1)
\end{cousinweb}

\begin{grammarlens}[title={What you've built}]
One more word for this chapter.
\end{grammarlens}

\subsection*{Your turn}
\begin{itemize}
\raggedright
  \item \emph{Say it:} \emph{la tour}
  \item \emph{Say it:} \emph{la tour}, once more
  \item \emph{Say it:} say \emph{la tour}
  \item \emph{Recall:} say the French for a guide, a guidebook, then the French for a castle, then say \emph{la tour} again
\end{itemize}

\begin{tcolorbox}[breakable,colback=teal!4,colframe=teal!35!black,title={Before you move on}]
What does la tour mean? (\textquotedblleft{}A tower\textquotedblright{}.) Say it once more.
\end{tcolorbox}
\section[le monument]{le monument — a monument}

\label{lesson:FR-C170-monument}



\begin{quote}
Before the new one: say the French for a castle, then the French for a tower.
\end{quote}

\subsection*{You'll want to know: le monument}
\textbf{le monument} — \textquotedblleft{}a monument\textquotedblright{}.

\begin{cousinweb}
\textbf{monument} $\leftarrow$ Latin \emph{monumentum}, \textquotedblleft{}a memorial\textquotedblright{}, from \emph{monere}, \textquotedblleft{}to remind\textquotedblright{} — a learned borrowing. (EtymDB 2.1)
\end{cousinweb}

\begin{grammarlens}[title={What you've built}]
One more word for this chapter.
\end{grammarlens}

\subsection*{Your turn}
\begin{itemize}
\raggedright
  \item \emph{Say it:} \emph{le monument}
  \item \emph{Say it:} \emph{le monument}, once more
  \item \emph{Say it:} say \emph{le monument}
  \item \emph{Recall:} say the French for a castle, then the French for a tower, then say \emph{le monument} again
\end{itemize}

\begin{tcolorbox}[breakable,colback=teal!4,colframe=teal!35!black,title={Before you move on}]
What does le monument mean? (\textquotedblleft{}A monument\textquotedblright{}.) Say it once more.
\end{tcolorbox}
\section[la statue]{la statue — a statue}

\label{lesson:FR-C170-statue}



\begin{quote}
Before the new one: say the French for a tower, then the French for a monument.
\end{quote}

\subsection*{You'll want to know: la statue}
\textbf{la statue} — \textquotedblleft{}a statue\textquotedblright{}.

\begin{cousinweb}
\textbf{statue} $\leftarrow$ Latin \emph{statua}, \textquotedblleft{}a standing image\textquotedblright{}, from \emph{stare}, \textquotedblleft{}to stand\textquotedblright{} (Lewis \& Short) — a learned borrowing.
\end{cousinweb}

\begin{grammarlens}[title={What you've built}]
One more word for this chapter.
\end{grammarlens}

\subsection*{Your turn}
\begin{itemize}
\raggedright
  \item \emph{Say it:} \emph{la statue}
  \item \emph{Say it:} \emph{la statue}, once more
  \item \emph{Say it:} say \emph{la statue}
  \item \emph{Recall:} say the French for a tower, then the French for a monument, then say \emph{la statue} again
\end{itemize}

\begin{tcolorbox}[breakable,colback=teal!4,colframe=teal!35!black,title={Before you move on}]
What does la statue mean? (\textquotedblleft{}A statue\textquotedblright{}.) Say it once more.
\end{tcolorbox}
\section[la fontaine]{la fontaine — a fountain}

\label{lesson:FR-C170-fontaine}



\begin{quote}
Before the new one: say the French for a monument, then the French for a statue.
\end{quote}

\subsection*{You'll want to know: la fontaine}
\textbf{la fontaine} — \textquotedblleft{}a fountain\textquotedblright{}.

\begin{cousinweb}
\textbf{fontaine} $\leftarrow$ Late Latin \emph{fontana}, from \emph{fons}, \emph{fontis}, \textquotedblleft{}a spring\textquotedblright{} — inherited. (EtymDB 2.1)
\end{cousinweb}

\begin{grammarlens}[title={What you've built}]
One more word for this chapter.
\end{grammarlens}

\subsection*{Your turn}
\begin{itemize}
\raggedright
  \item \emph{Say it:} \emph{la fontaine}
  \item \emph{Say it:} \emph{la fontaine}, once more
  \item \emph{Say it:} say \emph{la fontaine}
  \item \emph{Recall:} say the French for a monument, then the French for a statue, then say \emph{la fontaine} again
\end{itemize}

\begin{tcolorbox}[breakable,colback=teal!4,colframe=teal!35!black,title={Before you move on}]
What does la fontaine mean? (\textquotedblleft{}A fountain\textquotedblright{}.) Say it once more.
\end{tcolorbox}
